\section{Weekly Summary}

\begin{tcolorbox}[colback=yellow!10,colframe=black!50,title=AI-Generated Content]
\noindent The summary below was written by Claude (Anthropic) from that week's regression outputs and series descriptions. It has not been reviewed or fact-checked by a human, and should be read as an automated, informal recap -- not analysis.
\end{tcolorbox}

This week, the trusty random-series-mashing machine reached deep into the FRED grab bag and pulled out two pairings that no sane economist would ever put in the same sentence, let alone the same regression. Naturally, we regressed them anyway, because that is the entire point of this deeply respectable scientific endeavor.

Monday's contestant was a face-off between the Labor Compensation for Electric Power Utilities and, of all things, the Economic Policy Uncertainty Index for South Korea (a series so committed to the bit that it has been discontinued). And wouldn't you know it, the p-value came in at a positively swaggering 0.0006, with an R-squared of 0.41. Hold the phone. Forty-one percent of the variation, and a slope of 2.28 that says every uptick in Korean policy jitters comes with a hearty bump in American electrical worker paychecks. Is this the hidden thread binding the global economy together? Have we cracked it? No. We have not. It is two unrelated annual-ish time series that happened to drift in vaguely the same direction over a few decades, which is roughly as meaningful as noticing that your houseplant grew while your rent went up.

Tuesday somehow topped it. We paired Total Net Loan Charge-offs for tiny banks in the Middle Atlantic (also discontinued, because apparently this project only dates series that have already left the building) against Home Price Sales Pair Counts for Tampa, Florida. The relationship was gloriously negative: a slope of nearly negative 4,913, meaning that as those small-bank loans went sour, Tampa home sales allegedly went absolutely feral. The R-squared was a more humble 0.27, but the p-value strolled in at ten-to-the-minus-ten, which is the statistical equivalent of the machine slamming its fist on the table and yelling "TRUST ME." Please do not trust it.

As always, the mandatory reminder: these are randomly paired economic series with zero causal story, zero theoretical justification, and zero business being interpreted as anything other than a party trick. Spurious correlation is not the occasional embarrassing accident here -- it is the house specialty, served fresh daily. Small p-values feel dramatic, but with enough random dice rolls you will eventually roll a Yahtzee that means nothing. So enjoy the pretty lines, admire the confident tables, and under no circumstances rebalance your portfolio based on Tampa real estate and the mood of South Korean policymakers. See you next week for more numbers pretending to be friends.
